\begin{papercard}{benjumea2020}
\cf{Stages} \stagetag{A} \stagetag{D} \stagetag{C} \stagetag{H}. A scoping review of \emph{how} mHealth app privacy has been assessed: components, criteria, sources of criteria and scoring methods. It proposes a taxonomy of assessment objects and a list of 21 privacy items.

\cf{Aim} Understand how privacy is assessed for mHealth apps (components, scales, criteria, scoring) and propose a simple taxonomy \pc{benjumea2020}{1}{This scoping review aims to understand how privacy is assessed for mHealth apps}.

\cf{Review method} Arksey and O'Malley framework, reported with PRISMA-ScR \pc{benjumea2020}{2}{we conducted a scoping review using Arksey and O}. Search in July 2019 of Scopus, PubMed, IEEE Xplore and ACM DL with the string \emph{privacy AND (health app* OR mobile health OR mhealth) AND (test* OR framework OR review* OR evaluat* OR assess* OR ``comparative analysis'' OR ``regulation compliance'' OR taxonomy)} \pc{benjumea2020}{2}{The search string used was as follows}. Inclusion: papers that assess the privacy of mHealth apps, English title and abstract, 2009 onward. Excluded: papers without an assessment method (concerns, threats, challenges only) and papers evaluating no app \pc{benjumea2020}{2}{Papers that did not propose a method to evaluate privacy were}.

\cf{Screening and reliability} 710 records, minus 230 duplicates, gave 480 screened. Two authors independently screened 10\,\% (48) of titles and abstracts to check agreement: Cohen's $\kappa=0.73$. Each then screened half of the rest. 77 went to full text and 24 were included, with discrepancies resolved by consensus \pc{benjumea2020}{3}{Initially, two authors}. Data charting: an iteratively defined charting table with four groups (general information, evaluation procedure, evaluation criteria, scoring method), filled independently by two authors in Excel, with discrepancies resolved by consensus \pc{benjumea2020}{3}{The authors followed a collaborative and iterative process to}, \pc{benjumea2020}{4}{independently extracted data from the 24 selected articles}.

\cf{Taxonomy of assessment objects} Defined by consensus after full-text review. (1) App properties and behaviour: the app is used and data entered, e.g.\ login type or registration. (2) In-app information. (3) Personal information types. (4) App communications (traffic interception). (5) Static/dynamic analysis. (6) Existence of a privacy policy. (7) Content of the policy and/or terms. (8) Policy legibility (length, phrases, readability formulas) \pc{benjumea2020}{3}{App properties and behavior: this category refers to the app}.

\cf{21 privacy items} From policy content: data controller, DPO, purposes, legal basis, recipients, international transfers, data rights (incl.\ withdrawal of consent), obligation to provide data, profiling, nature of data, risks of collection, location of data, anonymisation. From in-app information: purposes, legal basis, recipients, rights, risks, protection of minors. Others: registration required, minimum data, policy existence, good practices \pc{benjumea2020}{3}{following items can be defined according to the content of a}. Only four items were used by more than half the studies: recipients identified (16/24), policy existence (13), purposes stated (13), data rights (13, mostly only access/control) \pc{benjumea2020}{12}{We have defined 21 items, but only}.

\cf{Findings on methods} 19/24 studies used the privacy policy. The criteria came from the literature (10), the authors' own judgement (9), and only rarely from law (HIPAA 3, GDPR 2, and none of the latter checked GDPR compliance) \pc{benjumea2020}{8}{Not many of the papers used legal frameworks or}. Scoring: 13 of 14 multi-item studies produced a score. Items were binary (6) or binary with intermediate values: 0/0.5/1, 0/1/2, $-1$/0/1 (7). Only two weighted the items \pc{benjumea2020}{12}{The items were assessed in a binary manner in 6 out of the 14 papers}, \pc{benjumea2020}{13}{proposed a weighted score and 1 article}. Sample sizes ranged from 10 to 600 apps, with a mode of 20 and a median below 51 \pc{benjumea2020}{6}{the number of apps analyzed in each article is disparate}. Table 3 lists every study's criteria and scoring rule, e.g.\ Scott/Adhikari risk (0--3) and safety (0--6) scores, Brüggemann's weighted privacy risk score, Zapata's 0/0.5/1 items \pc{benjumea2020}{9}{Criteria for evaluation of the apps}.

\cf{Critique} No study explained how it evaluated policy content for items such as secondary use. Hutton's raters disagreed over unclear policy language, but no one specified exact criteria, which threatens reproducibility \pc{benjumea2020}{14}{None of the articles explained how they evaluated}. Studies looked in different places: both policy and app (2), app only (6), unclear (4), policy and terms only (3). Searching in the wrong place can miss information \pc{benjumea2020}{14}{We consider that it is possible to miss}. Recommendation: objective, clearly defined items, ideally with a legal basis and explicit application rules \pc{benjumea2020}{15}{A positive step forward would be the creation of a scale}.

\cf{Limitations stated} English titles only, four databases, keyword choice, no pre-existing taxonomy, so charting is subjective \pc{benjumea2020}{15}{Specifically, for this review, the absence of an existing}.

\cf{Use for this thesis} (1) A literature map and a way to find comparable instruments (Table 3, Table 4). (2) A justification for writing \emph{explicit decision rules} and specifying \emph{where} each code is looked for (policy, store, UI), since this is the main reproducibility gap identified. (3) A precedent for reporting the screening $\kappa$ on a 10\,\% double-screened sample. (4) Sample-size context: 20 apps is the mode, not a rule. It is not itself an instrument.
\end{papercard}
